%%%%%%%%%%%%%%%%%
% CV tailored for Thales - Ingénieur Sûreté de Fonctionnement et Safety
% Form inspired by the user's Intheair CV structure: clear context/objective blocks,
% concise bullets, strong role positioning.
%%%%%%%%%%%%%%%%%

\documentclass[8pt,a4paper,ragged2e,withhyper]{altacv}

\geometry{left=1.25cm,right=1.25cm,top=.5cm,bottom=1.cm,columnsep=1.2cm}

\usepackage{paracol}
\usepackage{fontawesome5}
\usepackage{hyperref}

\iftutex 
  \setmainfont{Roboto Slab}
  \setsansfont{Lato}
  \renewcommand{\familydefault}{\sfdefault}
\else
  \usepackage[rm]{roboto}
  \usepackage[defaultsans]{lato}
  \renewcommand{\familydefault}{\sfdefault}
\fi

\definecolor{SlateGrey}{HTML}{2E2E2E}
\definecolor{LightGrey}{HTML}{666666}
\definecolor{DarkPastelRed}{HTML}{8AC0D9}
\definecolor{PastelRed}{HTML}{00B0DA}
\definecolor{GoldenEarth}{HTML}{B4AD0C}

\colorlet{name}{black}
\colorlet{tagline}{PastelRed}
\colorlet{heading}{DarkPastelRed}
\colorlet{headingrule}{GoldenEarth}
\colorlet{subheading}{PastelRed}
\colorlet{accent}{PastelRed}
\colorlet{emphasis}{SlateGrey}
\colorlet{body}{LightGrey}

\definecolor{violet}{HTML}{5A2582}
\definecolor{rouge}{HTML}{F53B3B}
\definecolor{noir}{HTML}{00232C}
\definecolor{bleu}{HTML}{00B0DA}
\definecolor{rose}{HTML}{F831A0}
\definecolor{jaune}{HTML}{FFEC00}
\definecolor{vert}{HTML}{34AD6C}

\renewcommand{\namefont}{\Huge\rmfamily\bfseries}
\renewcommand{\personalinfofont}{\footnotesize}
\renewcommand{\cvsectionfont}{\LARGE\rmfamily\bfseries}
\renewcommand{\cvsubsectionfont}{\large\bfseries}
\renewcommand{\cvItemMarker}{{\small\textbullet}}
\renewcommand{\cvRatingMarker}{\faCircle}

\input{../../_shared/pubs-num.tex}
\addbibresource{../../_shared/sample.bib}

\begin{document}

\name{BOILEAU Guillaume}
\tagline{Ingénieur RAMS / Safety | Spatial, maîtrise des risques, exigences, fiabilité \& validation système}

\photoL{2.8cm}{../../_shared/moi}

\personalinfo{%
  \email{guillaume.boileaupro@gmail.com}
  \phone{+33 6 59 94 85 84}
  \location{Alpes-Maritimes, FRANCE}
  \linkedin{guillaume-boileau-891b81265}
  \researchgate{Guillaume-Boileau-2}
  \orcid{0000-0002-3576-6968}
}

\makecvheader
\vspace{-0.3cm}

\AtBeginEnvironment{itemize}{\small}
\columnratio{0.64}

\begin{paracol}{2}

\cvsection{Experience}

\cvevent{\textbf{Ingénieur intégration logiciel -- Interfaces terrain, support opérationnel \& analyse d’incidents}}{VEV Platform Services France}{2025 -- 2026}{Valbonne, France}

\textbf{Contexte :} Environnement industriel logiciel pour une plateforme CPMS liée à des infrastructures de recharge de véhicules électriques, avec services backend, interfaces applicatives, données opérationnelles et équipements terrain.

\textbf{Objectif :} Contribuer au développement, à l’intégration et au suivi technique de services connectés à des équipements terrain, en assurant la cohérence entre flux de données, comportements applicatifs, contraintes hardware/software et attentes opérationnelles.

\begin{itemize}
  \item \textbf{Analyse d’incidents :} Investigation de comportements anormaux, interruptions d’échanges, incohérences de données et effets sur le service.
  \item \textbf{Maîtrise des risques opérationnels :} Identification de points de fragilité dans les flux et interfaces, formalisation des constats et contribution aux actions de fiabilisation.
  \item \textbf{Intégration HW/SW :} Travail sur les interfaces entre services logiciels, données opérationnelles, infrastructures connectées et comportements terrain.
  \item \textbf{Validation technique :} Vérification de la cohérence entre données applicatives, comportement des équipements et attendus fonctionnels.
  \item \textbf{Développement logiciel :} Contribution à des composants TypeScript, Angular et Node.js intégrés dans une plateforme opérationnelle.
  \item \textbf{Coordination agile :} Suivi des tâches, priorisation, backlog, points d’avancement et coordination via Zenhub.
  \item \textbf{Reporting technique :} Synthèse d’anomalies, impacts, points bloquants et actions de correction pour faciliter les arbitrages.
\end{itemize}

\divider

\cvevent{\textbf{Ingénieur R\&D senior -- Systèmes spatiaux, risques, performance \& validation}}{CNES}{2023 -- 2025}{Nice, France}

\textbf{Contexte :} Activités d’ingénierie et de R\&D pour la mission spatiale LISA ESA/NASA, dans un environnement international impliquant simulation, développement Python, traitement de données, intégration de modèles, validation technique et contraintes de performance.

\textbf{Objectif :} Structurer, intégrer, comparer et valider des contributions techniques complexes afin de démontrer la cohérence des résultats, identifier les écarts critiques et produire des éléments exploitables pour revue et décision.

\begin{itemize}
  \item \textbf{Analyse de risques techniques :} Identification d’écarts entre modèles, hypothèses, sorties de simulation et résultats attendus ; analyse d’impact et proposition d’actions correctives.
  \item \textbf{Validation système :} Définition de contrôles de cohérence, règles de comparaison, méthodes de vérification et critères d’acceptabilité des résultats.
  \item \textbf{Exigences et traçabilité :} Structuration de workflows, configurations, versions, données d’entrée, sorties de simulation et résultats d’analyse pour assurer la traçabilité.
  \item \textbf{Performance et robustesse :} Analyse de simulations dans des contextes de signaux faibles, bruit, incertitudes, limites de performance et contraintes de cohérence.
  \item \textbf{Criticités techniques :} Remontée des points bloquants, anomalies de cohérence, risques résiduels et recommandations pour sécuriser les analyses.
  \item \textbf{Coordination technique :} Interface avec des contributeurs scientifiques, logiciels et institutionnels dans un environnement spatial international.
  \item \textbf{Préparation de revues :} Consolidation de livrables, synthèses techniques, traçabilité des hypothèses et préparation d’éléments de décision.
  \item \textbf{Plateforme logicielle :} Développement de CATpyLISA, plateforme Python pour l’intégration, la comparaison, la validation et la revue technique de modèles.
  \textcolor{DarkPastelRed}{\href{https://gitlab.oca.eu/gboileau/lisa_l3_tools}{\bf GitLab: CATpyLISA}}
\end{itemize}

\divider

\cvevent{\textbf{Data Scientist -- Risques de bruit, modélisation \& analyse de performance}}{Universiteit Antwerpen, Belgium}{2021 -- 2023}{Antwerp, Belgium}

\textbf{Contexte :} Études de performance pour instruments de précision et détecteurs de nouvelle génération, combinant mesures environnementales, bruit, modélisation statistique, configurations instrumentales et validation de résultats.

\textbf{Objectif :} Développer des workflows robustes et reproductibles pour identifier, caractériser et réduire des sources affectant la qualité de mesure, la sensibilité et la disponibilité scientifique de systèmes complexes.

\begin{itemize}
  \item \textbf{Analyse de performance :} Évaluation de la qualité de mesure, des limites instrumentales et de l’impact de différentes sources de bruit.
  \item \textbf{Maîtrise des risques physiques :} Analyse de l’impact du bruit sismique et environnemental sur la sensibilité instrumentale et les budgets de bruit.
  \item \textbf{Données capteurs :} Analyse de mesures issues de capteurs environnementaux, corrélations spatiales et effets de propagation.
  \item \textbf{Modélisation statistique :} Développement de méthodes d’analyse et d’inférence bayésienne, notamment par chaînes MCMC.
  \item \textbf{Configurations système :} Études de différentes configurations instrumentales et de différents niveaux de bruit pour la détection de signaux très ténus.
  \item \textbf{Validation de résultats :} Analyse de sensibilité, vérification de cohérence, figures de mérite et évaluation des limites de performance.
  \item \textbf{Reporting technique :} Production de résultats traçables, de synthèses techniques et de recommandations pour parties prenantes internationales.
\end{itemize}

\divider

\cvevent{\textbf{Ingénieur R\&D junior -- Simulation, signal \& diagnostic de risques techniques}}{ARTEMIS, Observatoire de la Côte d’Azur}{2018 -- 2021}{Nice, France}

\textbf{Contexte :} Activités de R\&D liées à l’analyse de signaux faibles, à la simulation numérique, aux diagnostics de bruit et à la préparation de missions spatiales et d’instruments de détection.

\textbf{Objectif :} Développer des méthodes d’analyse, automatiser des simulations et produire des résultats exploitables dans un environnement scientifique et technique exigeant.

\begin{itemize}
  \item \textbf{Simulation numérique :} Développement d’approches de simulation pour environnements de signaux complexes et jeux de données de grande taille.
  \item \textbf{Analyse de données :} Mise en place de méthodes pour analyser et séparer des signaux faibles et superposés.
  \item \textbf{Diagnostic technique :} Analyse de sources de bruit, effets de corrélation entre capteurs et risques d’interprétation des résultats.
  \item \textbf{Validation technique :} Vérification de cohérence, comparaison de résultats et études de sensibilité.
  \item \textbf{Performance système :} Études de robustesse, incertitudes et limites de performance dans des environnements bruités.
  \item \textbf{Documentation :} Formalisation des méthodes, hypothèses, résultats et conclusions pour revue collaborative.
\end{itemize}

\switchcolumn

\cvsection{Profile}

\textbf{Ingénieur docteur orienté RAMS, Safety et validation système, avec une expérience solide en environnement spatial international, maîtrise des risques techniques, analyse de performance, simulation, traitement du signal, données capteurs et traçabilité des résultats.}

Mon parcours combine analyse de systèmes complexes, identification de criticités, validation de résultats, consolidation d’exigences techniques, recommandations de réduction de risques, documentation de revue et coordination transverse. J’ai travaillé sur des environnements exigeants associant logiciel, données, instrumentation, bruit, performance, fiabilité des analyses et contraintes opérationnelles.

\begin{itemize}
  \item Expérience en analyse de risques techniques, validation système, exigences, traçabilité, performances et démonstration de conformité des résultats.
  \item Culture forte du spatial, des systèmes complexes, des interfaces multiples, de la simulation et des contraintes de sûreté/performance.
  \item Capacité à identifier les criticités, remonter les risques résiduels et formuler des recommandations exploitables.
  \item Profil rigoureux, organisé, autonome, à l’aise dans les environnements transnationaux et les interfaces avec métiers d’ingénierie.
\end{itemize}

\cvsection{Languages}

\cvskill{English}{4}
\divider
\cvskill{Spanish}{2}
\divider

\medskip

\cvsection{Education}

\cvevent{PhD in Physics}{Université Côte d’Azur}{2018 -- 2021}{Nice, France}

\divider

\cvevent{Selective Graduate Program in Fundamental Physics (Magistère)}{Université Paris-Saclay}{2016 -- 2018}{Orsay, France}

\divider

\cvevent{MSc in Astronomy and Astrophysics}{Observatoire de Paris / Université Paris-Saclay}{2017 -- 2018}{Paris, France}

\divider

\cvevent{BSc in Fundamental Physics}{Université Paris-Saclay}{2015 -- 2016}{Orsay, France}

\cvsection{Professional Training}

\cvevent{Machine Learning in Python with scikit-learn}{FUN MOOC -- Inria}{2026}{France Université Numérique}

\small{
Predictive modelling, supervised learning, model evaluation, cross-validation, metrics, pipelines and practical machine learning workflows with Python and scikit-learn.
}

\divider

\cvevent{Supervised Machine Learning: Regression \& Classification}{DeepLearning.AI -- Andrew Ng}{2020}{Coursera}

\divider

\cvevent{Recherche reproductible : principes méthodologiques pour une science transparente}{FUN MOOC -- Inria}{2025}{France Université Numérique}

\divider

\cvevent{Continuous Professional Training -- Software, Data, AI and Systems}{OpenClassrooms}{2024 -- 2026}{}

\small{
Python, SQL, Machine Learning, AI, Linux, JavaScript, TypeScript, Angular, Node.js, Django, REST APIs, C/C++, web development, algorithms and software engineering fundamentals.
}

\end{paracol}

\cvsection{Projects}

\cvevent{CATpyLISA -- Plateforme Python d’intégration, comparaison et validation}{Mission spatiale LISA ESA/NASA}{2023 -- 2025}{}

Développement d’une plateforme Python dédiée à l’intégration de modèles, à la structuration de données, à la comparaison de résultats, à la validation technique et à la revue de workflows liés à la mission LISA.

\textbf{Rôle :} développement Python, conception de pipelines, structuration de données, logique de validation, contrôles de cohérence, analyse d’écarts, documentation technique, consolidation de contributions et coordination avec plusieurs contributeurs.

\textbf{Compétences mobilisées :} Python, intégration de modèles, validation système, analyse de risques techniques, traitement du signal, analyse de performance, simulation, documentation, GitLab, Confluence, environnement spatial international.

\divider

\cvevent{Noise budget, performance and risk analysis workflows}{Instrumentation scientifique et systèmes de précision}{2021 -- 2023}{}

Développement de workflows numériques et statistiques pour analyser des contributions de bruit, comparer des configurations, évaluer la robustesse des performances et formuler des recommandations sur les budgets de bruit et configurations système.

\textbf{Rôle :} analyse de performance, bruit, simulation, traitement du signal, caractérisation de risques, études de sensibilité, diagnostics techniques, reporting et support à la décision.

\textbf{Compétences mobilisées :} traitement du signal, simulation numérique, analyse spectrale, Python, MCMC, validation, analyse de sensibilité, documentation technique.

\divider

\cvevent{VEV -- Interfaces hardware/software, incidents \& fiabilisation opérationnelle}{Industrial software / Connected infrastructure}{2025 -- 2026}{}

Contribution au développement, à l’intégration, au suivi opérationnel et à la validation de services logiciels connectés à des infrastructures terrain.

\textbf{Rôle :} développement logiciel, investigation d’incidents, analyse de flux de données, support opérationnel, suivi de tickets, coordination agile et reporting technique.

\textbf{Compétences mobilisées :} TypeScript, Angular, Node.js, FTP, Linux, Git, Zenhub, analyse d’anomalies, support production, interfaces terrain.

\cvsection{Compétences clés}

\vspace{-.3cm}

\cvsubsection{RAMS, Safety \& maîtrise des risques}

{\small
\cvtag{RAMS -- ramp-up}
\cvtag{Sûreté de fonctionnement}
\cvtag{Safety}
\cvtag{Maîtrise des risques}
\cvtag{Analyse de criticité}
\cvtag{Risques résiduels}
\cvtag{Recommandations techniques}
\cvtag{Fiabilité -- awareness}
\cvtag{Disponibilité -- awareness}
\cvtag{Tolérance aux défaillances -- awareness}
\cvtag{Démonstration de conformité}
\cvtag{Réglementations spatiales -- ramp-up}
\cvtag{SDM / STM -- ramp-up}
}

\divider\smallskip
\vspace{-.3cm}

\cvsubsection{Exigences, validation \& conformité}

{\small
\cvtag{Déclinaison d’exigences}
\cvtag{Vérification d’exigences}
\cvtag{Traçabilité}
\cvtag{IVV système}
\cvtag{Vérification / Validation}
\cvtag{Contrôles de cohérence}
\cvtag{Critères d’acceptation}
\cvtag{Préparation de revues}
\cvtag{Documentation qualité}
\cvtag{Analyse d’écarts}
\cvtag{Root cause analysis}
\cvtag{Diagnostic technique}
\cvtag{Confluence}
}

\divider\smallskip
\vspace{-.3cm}

\cvsubsection{Systèmes complexes, spatial \& ingénierie}

{\small
\cvtag{Spatial}
\cvtag{Systèmes satellitaires}
\cvtag{LISA ESA/NASA}
\cvtag{CNES}
\cvtag{ESA}
\cvtag{Instrumentation scientifique}
\cvtag{Systèmes complexes}
\cvtag{Modélisation système}
\cvtag{Simulation numérique}
\cvtag{EGSE -- ramp-up}
\cvtag{Avant-projets -- awareness}
\cvtag{Interfaces métiers}
\cvtag{Environnement transnational}
\cvtag{Anglais courant}
}

\divider\smallskip
\vspace{-.3cm}

\cvsubsection{Performance, signal \& données}

{\small
\cvtag{Analyse de performance}
\cvtag{Traitement du signal}
\cvtag{Analyse de bruit}
\cvtag{Budget de bruit}
\cvtag{Signaux faibles}
\cvtag{Low-SNR}
\cvtag{Analyse de sensibilité}
\cvtag{Comparaison de modèles}
\cvtag{Data quality}
\cvtag{Données capteurs}
\cvtag{Statistiques}
\cvtag{MCMC}
\cvtag{Machine Learning}
}

\divider\smallskip
\vspace{-.3cm}

\cvsubsection{Logiciel, data \& automatisation}

{\small
\cvtag{Python}
\cvtag{C}
\cvtag{Pandas}
\cvtag{NumPy}
\cvtag{SciPy}
\cvtag{Matplotlib}
\cvtag{scikit-learn}
\cvtag{TypeScript}
\cvtag{Angular}
\cvtag{Node.js}
\cvtag{REST API}
\cvtag{SQL}
\cvtag{Bash}
\cvtag{Linux}
\cvtag{Git}
\cvtag{GitLab}
\cvtag{Docker}
\cvtag{CI/CD}
}

\divider\smallskip
\vspace{-.3cm}

\cvsubsection{Coordination, reporting \& outils}

{\small
\cvtag{Coordination technique}
\cvtag{Interfaces client / métiers}
\cvtag{Support décisionnel}
\cvtag{Reporting}
\cvtag{KPIs}
\cvtag{Synthèse technique}
\cvtag{Organisation}
\cvtag{Qualités relationnelles}
\cvtag{Jira}
\cvtag{Zenhub}
\cvtag{Power BI}
\cvtag{OpenSearch}
\cvtag{Sentry}
\cvtag{Autonomie}
}

\end{document}
